\documentclass[11pt]{beamer}

% ============================================================
% Theme
% ============================================================
\usetheme{Madrid}
\usecolortheme{seahorse}
\setbeamertemplate{navigation symbols}{}

% ============================================================
% Encoding
% ============================================================
\usepackage[utf8]{inputenc}
\usepackage[T1]{fontenc}
\usepackage{lmodern}

% ============================================================
% Packages
% ============================================================
\usepackage{amsmath}
\usepackage{listings}
\usepackage{xcolor}
\usepackage{tikz}

\lstset{
	basicstyle=\ttfamily\small,
	frame=single,
	breaklines=true,
	columns=fullflexible,
	keywordstyle=\color{blue},
	commentstyle=\color{gray},
	stringstyle=\color{red},
	showstringspaces=false
}

% ============================================================
% Title
% ============================================================
\title{Concurrency \& Distributed Systems}
\subtitle{Week 1 --- Part 4: Measurement-Driven Engineering: EnhancedLoadClient}
\author{Dr. Mohamad Aoude}
\institute{Lebanese University \\ Faculty of Engineering}
\date{\today}

\begin{document}
	
	% ============================================================
	\begin{frame}
		\titlepage
	\end{frame}
	
	% ============================================================
	\section{Why We Need a Real Load Client}
	
	% ------------------------------------------------------------
	\begin{frame}{Demo is Not Enough}
		Previously:
		\begin{itemize}
			\item We ran a simple LoadClient
			\item We observed “it feels slower”
			\item We guessed performance
		\end{itemize}
		

		
		\begin{alertblock}{Problem}
			Guessing is not engineering.
		\end{alertblock}

		
		\begin{block}{Goal}
			We now measure:
			\begin{itemize}
				\item p50 latency
				\item p95 latency
				\item Throughput (requests/sec)
			\end{itemize}
		\end{block}
	\end{frame}
	
	% ============================================================
	\section{What EnhancedLoadClient Measures}
	
	% ------------------------------------------------------------
	\begin{frame}{Latency Percentiles}
		\textbf{p50 (Median)}
		\begin{itemize}
			\item 50\% of users are faster than this
			\item Represents “typical experience”
		\end{itemize}
		

		
		\textbf{p95 (Tail)}
		\begin{itemize}
			\item 95\% of users are faster
			\item Slowest 5\% determine perceived quality
		\end{itemize}

		
		\begin{alertblock}{Critical Insight}
			High throughput with bad p95 = bad system.
		\end{alertblock}
	\end{frame}
	
	% ------------------------------------------------------------
	\begin{frame}{Throughput}
		\[
		\text{Throughput} =
		\frac{\text{Total Requests}}{\text{Total Time}}
		\]

		
		If:
		\begin{itemize}
			\item 500 requests completed
			\item in 2 seconds
		\end{itemize}
		
		\[
		Throughput = 250 \text{ req/sec}
		\]

		
		\begin{block}{Remember}
			Latency measures individual experience.  
			Throughput measures system capacity.
		\end{block}
	\end{frame}
	
	% ============================================================
	\section{How To Run It}
	
	% ------------------------------------------------------------
	\begin{frame}[fragile]{Step 1: Start Server}
		Terminal 1:
		
		\begin{lstlisting}[language=bash]
			javac ImprovedMultiThreadedServer.java
			java ImprovedMultiThreadedServer
		\end{lstlisting}
	\end{frame}
	
	% ------------------------------------------------------------
	\begin{frame}[fragile]{Step 2: Compile Client}
		Terminal 2:
		
		\begin{lstlisting}[language=bash]
			javac EnhancedLoadClient.java
		\end{lstlisting}
	\end{frame}
	
	% ------------------------------------------------------------
	\begin{frame}[fragile]{Step 3: Run Experiments}
		
		Light load:
		
		\begin{lstlisting}[language=bash]
			java EnhancedLoadClient --clients=10 --requests=10
		\end{lstlisting}

		
		Medium load:
		
		\begin{lstlisting}[language=bash]
			java EnhancedLoadClient --clients=50 --requests=20
		\end{lstlisting}

		
		Heavy load:
		
		\begin{lstlisting}[language=bash]
			java EnhancedLoadClient --clients=200 --requests=20
		\end{lstlisting}
	\end{frame}
	
	% ============================================================
	\section{What Students Must Observe}
	
	% ------------------------------------------------------------
	\begin{frame}{Single-Threaded Server}
		Expected behavior:
		
		\begin{itemize}
			\item Latency grows linearly
			\item p95 extremely high
			\item Throughput capped
			\item CPU mostly idle
		\end{itemize}

		
		\begin{alertblock}
			Queue formation dominates.
		\end{alertblock}
	\end{frame}
	
	% ------------------------------------------------------------
	\begin{frame}{Multi-Threaded Server}
		Expected behavior:
		
		\begin{itemize}
			\item Latency collapses initially
			\item Throughput increases significantly
			\item CPU rises moderately
		\end{itemize}

		
		Then push further:
		
		\begin{itemize}
			\item p95 increases again
			\item Throughput plateaus
			\item CPU spikes
		\end{itemize}
	\end{frame}
	
	% ============================================================
	\section{Sweet Spot Concept}
	
	% ------------------------------------------------------------
	\begin{frame}{Why Too Many Threads Hurt}
		
		\textbf{Three causes:}
		
		\begin{itemize}
			\item Context switching overhead
			\item Memory overhead (stack per thread)
			\item Lock contention
		\end{itemize}

		
		\[
		\text{Optimal Threads}
		\approx
		\text{cores} \times
		\left(1 + \frac{Wait}{Compute}\right)
		\]
	\end{frame}
	
	% ============================================================
	\section{Lab Instructions}
	
	% ------------------------------------------------------------
	\begin{frame}{Lab 1.1 Requirements}
		
		Students must submit:
		
		\begin{itemize}
			\item Measurement table (p50, p95, throughput)
			\item At least 3 pool sizes tested
			\item CSV export file
			\item 5--8 sentence engineering explanation
		\end{itemize}
	\end{frame}
	
	% ------------------------------------------------------------
	\begin{frame}{Required Measurement Table}
		
		\begin{center}
			\begin{tabular}{lccc}
				Load & p50 & p95 & Throughput \\
				\hline
				Light (10) &  &  &  \\
				Medium (50) &  &  &  \\
				Heavy (100) &  &  &  \\
			\end{tabular}
		\end{center}

		
		\begin{alertblock}
			No numbers = no credit.
		\end{alertblock}
	\end{frame}
	
	% ============================================================
	\section{Interpretation Questions}
	
	% ------------------------------------------------------------
	\begin{frame}{Students Must Answer}
		
		\begin{enumerate}
			\item Why did latency drop with threads?
			\item Why did p95 increase under heavy load?
			\item Why didn’t throughput grow infinitely?
			\item Was the system I/O-bound or CPU-bound?
		\end{enumerate}

		
		\begin{block}{Engineering Language Required}
			Use:
			\begin{itemize}
				\item queue
				\item overlap
				\item scheduling overhead
				\item saturation
			\end{itemize}
		\end{block}
	\end{frame}
	
	% ============================================================
	\section{Final Synthesis}
	
	% ------------------------------------------------------------
	\begin{frame}{What This Demonstrates}
		
		\begin{itemize}
			\item Concurrency hides waiting.
			\item Concurrency does NOT create CPU power.
			\item Tail latency reveals overload.
			\item Measurement drives architecture decisions.
		\end{itemize}

		
		\begin{alertblock}
			If you cannot measure it,  
			you cannot engineer it.
		\end{alertblock}
	\end{frame}
	
	% ------------------------------------------------------------
	\begin{frame}{Next Step}
		
		\Large
		What happens when threads share mutable state?

		
		\begin{alertblock}
			Next Lecture: Race Conditions, Visibility, and Deadlocks.
		\end{alertblock}
		
	\end{frame}
	
\end{document}
